\section*{Two-point overlap surgery: exact spectrum, entropy, influence, and upper-set potential}

All logarithms are base two.  Let $d\ge2$, let $z\in\mathbb F_2^d$ have
Hamming weight $h\in\{1,\ldots,d\}$, and define the genuine Boolean
function
\[
 f_{d,h}(x)=\begin{cases}-1,&x=0\text{ or }x=z,\\+1,&\text{otherwise}.
 \end{cases}                                                   \tag{1}
\]
Equivalently, $f_{d,h}$ is the product of two one-point defects on the
same physical $d$-cube.  Thus (1) is shared-input overlap, not a disjoint
tensor product.  Put $\delta=2^{-d}$ and
\[
 L_z=\{S\subseteq[d]:|S\cap\operatorname{supp}z|\equiv0\pmod2\}.
\]

\paragraph{Theorem.}
The squared Fourier spectrum of $f_{d,h}$ is
\[
 p_\varnothing=(1-4\delta)^2,\qquad
 p_S=16\delta^2\quad(S\in L_z\setminus\{\varnothing\}),       \tag{2}
\]
and is zero otherwise.  Consequently, with
$m=8\delta(1-2\delta)$ and $M=2^{d-1}-1$,
\[
 H(f_{d,h})=h_2(m)+m\log_2M.                                  \tag{3}
\]
The total influence is
\[
 I(f_{d,h})=\begin{cases}4(d-1)\delta,&h=1,\\4d\delta,&h\ge2.
 \end{cases}                                                   \tag{4}
\]

For a squared spectral law $p$, define the upper-set mass
$\pi(S)=\sum_{T\supseteq S}p_T$ and
\[
 U(p)=\sum_{S:p_S>0}p_S\log_2\frac{\pi(S)}{p_S}.
\]
Then
\[
 U(f_{d,1})=(1-4\delta)^2\log_2\frac1{(1-4\delta)^2}
 +4(d-1)\delta(1-4\delta),                                    \tag{5}
\]
while for $h\ge2$,
\[
 \begin{split}
 U(f_{d,h})={}&(1-4\delta)^2\log_2\frac1{(1-4\delta)^2}\\
 &+4(d-2)\delta-16(d-1)\delta^2
 +\mathbf1_{\{h\ {
m even}\}},2^{4-h}\delta.              \tag{6}
 \end{split}
\]
For every admissible $d,h$,
\[
 H(f_{d,h})<4I(f_{d,h}),\qquad U(f_{d,h})<3I(f_{d,h}).          \tag{7}
\]
If $h=h_d$ is arbitrary with $2\le h_d\le d$, then
\[
 \frac{H(f_{d,h_d})}{I(f_{d,h_d})}\longrightarrow4,
 \qquad
 \frac{U(f_{d,h_d})}{I(f_{d,h_d})}\longrightarrow1.          \tag{8}
\]

\paragraph{Proof.}
For $S\ne\varnothing$, Fourier expansion of the two flipped atoms gives
\[
 \widehat f(S)=-2\delta(1+\chi_S(z)).
\]
Together with $\widehat f(\varnothing)=1-4\delta$, this proves (2), and
(3) follows because the nonempty spectral mass is uniform on $M$ atoms.

The negative set has two vertices.  If $h=1$, they share one internal edge,
so their edge boundary has $2(d-1)$ undirected edges.  If $h\ge2$, it has
$2d$ edges.  Dividing by $2^{d-1}$ proves (4).

For nonempty $S\in L_z$, the number of $T\in L_z$ containing $S$ is
\[
 \#\{T\supseteq S:T\in L_z\}=
 \begin{cases}
 2^{d-|S|},&\operatorname{supp}z\subseteq S\text{ and }h\text{ is even},\\
 2^{d-|S|-1},&\operatorname{supp}z\not\subseteq S.
 \end{cases}                                                   \tag{9}
\]
When $h$ is odd the first case is impossible for $S\in L_z$.  Summing the
logarithm of (9) over the uniform nonempty spectrum gives (5)--(6).  For
$h=1$ this also reflects the exact identity of $f_{d,1}$ with a point defect
on $d-1$ effective coordinates.

The known point-defect inequalities give (7) when $h=1$.  For $h\ge2$,
$h_2(m)\le m\log_2(e/m)$ and $M<2^{d-1}$ yield
\[
 H\le8\delta(1-2\delta)
 \left[2d-4+\log_2e-\log_2(1-2\delta)\right]<16d\delta=4I,
\]
because $\delta\le1/4$ and
$\log_2e-\log_2(1-2\delta)<4$.  Also
\[
 (1-4\delta)^2\log_2\frac1{(1-4\delta)^2}
 \le\frac{8\delta(1-4\delta)}{\ln2},
\]
and the nonempty part of (6) is at most $4(d-1)\delta$.  Hence
$U<12d\delta=3I$ for $d\ge2$.  Finally, (3) and (6), uniformly in $h\ge2$,
give (8).\qed

\paragraph{Interpretation.}
The shared-input product destroys tensor factorization and, at $h=1$,
physically cancels one full coordinate's boundary contribution.  Nevertheless
both actual universal ratios remain bounded, with the same limits $4$ and
$1$ as a single rare defect.  Thus two-site overlap does not rescue the
divergent biased helper.  Any viable counterfamily must use a growing
nonlinear cluster geometry rather than finitely many overlapping atoms.
